\documentclass[10pt,a4paper]{amsart}
\usepackage{amsmath,amssymb,amsthm,graphicx}
% Retained minimal preamble for the Stacks Project chapter flat.tex (More on Flatness).
% Verbatim excerpt of the chapter preamble: Xy-pic setup (native preamble lines 14-18) and the native
% theorem declarations with the mathematical macros (native preamble lines 154-217). Copied mechanically;
% nothing added, removed or redefined. The wrapper class and the amsmath/amssymb/amsthm/graphicx base
% packages are supplied by the assembler, so the class/externaldocument part of the native preamble is omitted.
% For commutative diagrams we use Xy-pic
\usepackage[all]{xy}

% We use 2cell for 2-commutative diagrams.
\xyoption{2cell}
% Theorem environments.
%
\theoremstyle{plain}
\newtheorem{theorem}[subsection]{Theorem}
\newtheorem{proposition}[subsection]{Proposition}
\newtheorem{lemma}[subsection]{Lemma}

\theoremstyle{definition}
\newtheorem{definition}[subsection]{Definition}
\newtheorem{example}[subsection]{Example}
\newtheorem{exercise}[subsection]{Exercise}
\newtheorem{situation}[subsection]{Situation}

\theoremstyle{remark}
\newtheorem{remark}[subsection]{Remark}
\newtheorem{remarks}[subsection]{Remarks}

\numberwithin{equation}{subsection}

% Macros
%
\def\lim{\mathop{\mathrm{lim}}\nolimits}
\def\colim{\mathop{\mathrm{colim}}\nolimits}
\def\Spec{\mathop{\mathrm{Spec}}}
\def\Hom{\mathop{\mathrm{Hom}}\nolimits}
\def\Ext{\mathop{\mathrm{Ext}}\nolimits}
\def\SheafHom{\mathop{\mathcal{H}\!\mathit{om}}\nolimits}
\def\SheafExt{\mathop{\mathcal{E}\!\mathit{xt}}\nolimits}
\def\Sch{\mathit{Sch}}
\def\Mor{\mathop{\mathrm{Mor}}\nolimits}
\def\Ob{\mathop{\mathrm{Ob}}\nolimits}
\def\Sh{\mathop{\mathit{Sh}}\nolimits}
\def\NL{\mathop{N\!L}\nolimits}
\def\CH{\mathop{\mathrm{CH}}\nolimits}
\def\proetale{{pro\text{-}\acute{e}tale}}
\def\etale{{\acute{e}tale}}
\def\QCoh{\mathit{QCoh}}
\def\Ker{\mathop{\mathrm{Ker}}}
\def\Im{\mathop{\mathrm{Im}}}
\def\Coker{\mathop{\mathrm{Coker}}}
\def\Coim{\mathop{\mathrm{Coim}}}

% Boxtimes
%
\DeclareMathSymbol{\boxtimes}{\mathbin}{AMSa}{"02}

%
% Macros for moduli stacks/spaces
%
\def\QCohstack{\mathcal{QC}\!\mathit{oh}}
\def\Cohstack{\mathcal{C}\!\mathit{oh}}
\def\Spacesstack{\mathcal{S}\!\mathit{paces}}
\def\Quotfunctor{\mathrm{Quot}}
\def\Hilbfunctor{\mathrm{Hilb}}
\def\Curvesstack{\mathcal{C}\!\mathit{urves}}
\def\Polarizedstack{\mathcal{P}\!\mathit{olarized}}
\def\Complexesstack{\mathcal{C}\!\mathit{omplexes}}
% \Pic is the operator that assigns to X its picard group, usage \Pic(X)
% \Picardstack_{X/B} denotes the Picard stack of X over B
% \Picardfunctor_{X/B} denotes the Picard functor of X over B
\def\Pic{\mathop{\mathrm{Pic}}\nolimits}
\def\Picardstack{\mathcal{P}\!\mathit{ic}}
\def\Picardfunctor{\mathrm{Pic}}
\def\Deformationcategory{\mathcal{D}\!\mathit{ef}}
\usepackage{cleveref}
\crefname{originallemma}{original lemma}{original lemmas}
\Crefname{originallemma}{Original lemma}{Original lemmas}
\makeatletter
\newlabel{categories-lemma-right-localization}{{04VH}{source}{}{sourcepointer-1}{}}
\newlabel{categories-lemma-right-localization@cref}{{[originallemma][0][]04VH}{[1][0][]source}{}{}{}}
\makeatother
\crefname{originallemma}{original lemma}{original lemmas}
\Crefname{originallemma}{Original lemma}{Original lemmas}
\makeatletter
\newlabel{lemma-right-multiplicative-system}{{0ATV}{source}{}{sourcepointer-2}{}}
\newlabel{lemma-right-multiplicative-system@cref}{{[originallemma][0][]0ATV}{[1][0][]source}{}{}{}}
\makeatother
\crefname{originallemma}{original lemma}{original lemmas}
\Crefname{originallemma}{Original lemma}{Original lemmas}
\makeatletter
\newlabel{categories-lemma-properties-right-localization}{{04VK}{source}{}{sourcepointer-3}{}}
\newlabel{categories-lemma-properties-right-localization@cref}{{[originallemma][0][]04VK}{[1][0][]source}{}{}{}}
\makeatother
\crefname{originallemma}{original lemma}{original lemmas}
\Crefname{originallemma}{Original lemma}{Original lemmas}
\makeatletter
\newlabel{lemma-compactifyable}{{0A9Z}{source}{}{sourcepointer-4}{}}
\newlabel{lemma-compactifyable@cref}{{[originallemma][0][]0A9Z}{[1][0][]source}{}{}{}}
\makeatother
\crefname{originallemma}{original lemma}{original lemmas}
\Crefname{originallemma}{Original lemma}{Original lemmas}
\makeatletter
\newlabel{categories-lemma-morphisms-right-localization}{{04VI}{source}{}{sourcepointer-5}{}}
\newlabel{categories-lemma-morphisms-right-localization@cref}{{[originallemma][0][]04VI}{[1][0][]source}{}{}{}}
\makeatother
\title{Localizing the category of compactifications at the right multiplicative system recovers the category of compactifyable schemes}
\author{The Stacks Project Authors}
\date{stacks-project:flat.tex:local-snapshot-20261011; GNU FDL 1.2 or later}
\begin{document}
\maketitle
\noindent\textit{Editorial note.} All delivered mathematical text is the author's own, copied from the pinned Stacks Project chapter file flat.tex; no statement or proof step was rewritten, reordered or removed.\par\medskip
\noindent\textit{Editorial note.} The preamble is the chapter's own retained minimal preamble (sources/stacks-flat/retained-preamble.tex), a verbatim excerpt of the native Xy-pic setup and of the native theorem declarations and mathematical macros; no macro was invented or redefined, and the wrapper class and the amsmath/amssymb/amsthm/graphicx base packages are supplied by the assembler.\par\medskip
\noindent\textit{Editorial note.} The counted statement and proof contain no cross reference, no citation and no included file, so the delivered body is self contained.\par\medskip
\noindent\textit{Editorial note.} Canonical location of the counted proposition: Stacks tag 0ATW (sources/stacks-flat/pinned-original-source.tex line 9476).\par\medskip
\noindent\textit{Allowed standard prerequisites (editorial).}\begin{itemize}
\item Localization of a category at a right multiplicative system, its calculus of fractions and its functoriality. Categories, Lemmas labelled categories-lemma-right-localization (Stacks tag 04VH), categories-lemma-properties-right-localization (Stacks tag 04VK) and categories-lemma-morphisms-right-localization (Stacks tag 04VI), invoked at statement line 9479 and proof lines 9489, 9517. Supply the localized category C', the fact that the chosen multiplicative system is compatible with the functor, and the calculus of fractions that lets the two representations be assumed to share source and refinement; the equalizer construction itself is in the counted span.
\item Existence of compactifications, absolutely and relative to a proper base. Same chapter, Lemma labelled lemma-compactifyable (source line 9387) and Lemma labelled lemma-right-multiplicative-system (source line 9425), invoked at statement line 9482 and proof line 9500. Supply that a compactifyable scheme embeds into a proper scheme over S and that a morphism to a proper base may be compactified; the construction of the representing morphism of compactifications is inside the counted span.
\end{itemize}\medskip
\ifdefined\hypertarget\hypertarget{sourcepointer-1}{}\fi
\noindent\textit{Source reference.} Original lemma 04VH: sources/stacks-chow/tags.txt, source lines 6325--6325. Cross-chapter standard prerequisite located in the project tag list staged in the local snapshot; invoked at proof line None of the counted proof. It is retained only as a navigation pointer; every load-bearing step is inside the delivered span.\par\medskip
\ifdefined\hypertarget\hypertarget{sourcepointer-2}{}\fi
\noindent\textit{Source reference.} Original lemma 0ATV: sources/stacks-flat/pinned-original-source.tex, source lines 9425--9425. Same-chapter pointer into the author environment lemma at source line 9425; invoked at proof line None of the counted proof. It is retained only as a navigation pointer; every load-bearing step is inside the delivered span.\par\medskip
\ifdefined\hypertarget\hypertarget{sourcepointer-3}{}\fi
\noindent\textit{Source reference.} Original lemma 04VK: sources/stacks-chow/tags.txt, source lines 6328--6328. Cross-chapter standard prerequisite located in the project tag list staged in the local snapshot; invoked at proof line 9490 of the counted proof. It is retained only as a navigation pointer; every load-bearing step is inside the delivered span.\par\medskip
\ifdefined\hypertarget\hypertarget{sourcepointer-4}{}\fi
\noindent\textit{Source reference.} Original lemma 0A9Z: sources/stacks-flat/pinned-original-source.tex, source lines 9387--9387. Same-chapter pointer into the author environment lemma at source line 9387; invoked at proof line 9500 of the counted proof. It is retained only as a navigation pointer; every load-bearing step is inside the delivered span.\par\medskip
\ifdefined\hypertarget\hypertarget{sourcepointer-5}{}\fi
\noindent\textit{Source reference.} Original lemma 04VI: sources/stacks-chow/tags.txt, source lines 6326--6326. Cross-chapter standard prerequisite located in the project tag list staged in the local snapshot; invoked at proof line 9517 of the counted proof. It is retained only as a navigation pointer; every load-bearing step is inside the delivered span.\par\medskip
\begin{lemma}
\label{lemma-invert-right-multiplicative-system}
Let $S$ be a quasi-compact and quasi-separated scheme.
The functor $(X, \overline{X}) \mapsto X$ defines an
equivalence from the category of compactifications localized
(Categories, Lemma \ref{categories-lemma-right-localization})
at the right
multiplicative system of Lemma \ref{lemma-right-multiplicative-system}
to the category of compactifyable schemes over $S$.
\end{lemma}

\begin{proof}
Denote $\mathcal{C}$ the category of compactifications and
denote $Q : \mathcal{C} \to \mathcal{C}'$ the localization
functor of Categories, Lemma
\ref{categories-lemma-properties-right-localization}.
Denote $\mathcal{D}$ the category of compactifyable schemes
over $S$. It is clear from the lemma just cited and our
choice of multiplicative system that we
obtain a functor $\mathcal{C}' \to \mathcal{D}$.
This functor is clearly essentially surjective.
If $f : X \to Y$ is a morphism of compactifyable
schemes, then we choose an open immersion $Y \to \overline{Y}$
into a scheme proper over $S$, and then we choose an embedding
$X \to \overline{X}$ into a scheme $\overline{X}$ proper over
$\overline{Y}$ (possible by Lemma \ref{lemma-compactifyable}
applied to $X \to \overline{Y}$). This gives a morphism
$(X, \overline{X}) \to (Y, \overline{Y})$ of compactifications
which produces our given morphism $X \to Y$.
Finally, suppose given a pair of morphisms in the
localized category with the same source and target: say
$$
a = ((f, \overline{f}) : (X', \overline{X}') \to (Y, \overline{Y}),
(u, \overline{u}) : (X', \overline{X}') \to (X, \overline{X}))
$$
and
$$
b = ((g, \overline{g}) : (X'', \overline{X}'') \to (Y, \overline{Y}),
(v, \overline{v}) : (X'', \overline{X}'') \to (X, \overline{X}))
$$
which produce the same morphism $X \to Y$ over $S$, in other words
$f \circ u^{-1} = g \circ v^{-1}$. By
Categories, Lemma \ref{categories-lemma-morphisms-right-localization}
we may assume that $(X', \overline{X}') = (X'', \overline{X}'')$
and $(u, \overline{u}) = (v, \overline{v})$. In this case we
can consider the equalizer $\overline{X}''' \subset \overline{X}'$
of $\overline{f}$ and $\overline{g}$. The morphism
$(w, \overline{w}) : (X', \overline{X}''') \to (X', \overline{X}')$ is in
the multiplicative subset and we see that $a = b$ in the localized
category by precomposing with $(w, \overline{w})$.
\end{proof}

\end{document}
